\documentclass[11pt,a4paper]{article}
\usepackage[utf8]{inputenc}
\usepackage[spanish,es-nodecimaldot]{babel}
\usepackage{lmodern}
\usepackage{amsmath,amssymb,amsfonts}
\usepackage{geometry}
\geometry{margin=2.5cm}
\usepackage{titlesec}
\usepackage{enumitem}
\usepackage{booktabs}
\usepackage{array}
\usepackage{xcolor}
\usepackage{hyperref}

\titleformat{\section}{\large\bfseries\color{blue!70!black}}{\thesection}{1em}{}
\titleformat{\subsection}{\normalsize\bfseries\color{black}}{\thesubsection}{1em}{}
\titleformat{\subsubsection}{\small\bfseries\itshape\color{black!80}}{\thesubsubsection}{1em}{}

\title{\textbf{Facultad de Matemática, Astronomía, Física y Computación (FaMAF)}\\\large Didáctica Especial y Taller de Matemática\\\vspace{0.3cm}\normalsize Análisis de actividades relacionadas con la Gestión Curricular}
\author{\textbf{Resolución y Análisis Didáctico de Tareas (Excluyendo Tarea 7)}}
\date{\today}

\begin{document}

\maketitle

\section{Marco Teórico y Criterios de Análisis}

Para el desarrollo del presente trabajo se consideran tres ejes analíticos y didácticos fundamentales:
\begin{enumerate}
    \item \textbf{Gestión Curricular según Ponte (2005):} Clasificación de las tareas matemáticas según su \textit{grado de desafío} (reducido/alto) y su \textit{grado de apertura} (cerradas/abiertas):
    \begin{itemize}
        \item \textbf{Ejercicios:} Grado de desafío reducido, estructura cerrada (aplicación directa de algoritmos o conceptos conocidos).
        \item \textbf{Problemas:} Grado de desafío alto, estructura cerrada (requieren heurística, diseño de estrategias, pero la meta está claramente delimitada).
        \item \textbf{Actividades de Exploración:} Grado de desafío reducido/medio, estructura abierta (búsqueda de patrones, formulación de conjeturas iniciales).
        \item \textbf{Actividades de Investigación:} Grado de desafío alto, estructura abierta (los estudiantes definen preguntas, experimentan, conjeturan, prueban y comunican resultados).
    \end{itemize}
    \item \textbf{Ambientes de Aprendizaje según Skovsmose (2000):} Cruce entre el \textit{tipo de referencia} (Matemática Pura, Semirrealidad, Realidad) y la \textit{organización del trabajo} (Paradigma del Ejercicio, Escenario de Investigación), conformando seis ambientes:
    \begin{itemize}
        \item \textbf{Ambiente 1:} Matemática Pura / Paradigma del Ejercicio.
        \item \textbf{Ambiente 2:} Matemática Pura / Escenario de Investigación.
        \item \textbf{Ambiente 3:} Semirrealidad / Paradigma del Ejercicio.
        \item \textbf{Ambiente 4:} Semirrealidad / Escenario de Investigación.
        \item \textbf{Ambiente 5:} Realidad / Paradigma del Ejercicio.
        \item \textbf{Ambiente 6:} Realidad / Escenario de Investigación.
    \end{itemize}
    \item \textbf{Diseño Curricular de la Provincia de Córdoba:} Identificación del ciclo, año/curso, eje temático (Número y Operaciones, Álgebra y Funciones, Geometría y Medida, Estadística y Probabilidad) y los objetivos de aprendizaje correspondientes.
\end{enumerate}

\newpage

\section{Análisis y Resolución de las Tareas}

% ==========================================
% TAREA 1
% ==========================================
\subsection{Tarea 1}

\subsubsection{Enunciado y Desarrollo Matemático}
\begin{itemize}
    \item[a)] \textit{¿Es cierto que, sin hacer cuentas, se puede saber si $-3\cdot 15\cdot 14\cdot (-4) - 3\cdot 15\cdot 14\cdot (-4)$ da como resultado un número que es múltiplo de $5$ y de $35$?}
    
    \textbf{Resolución:}
    Observemos la operación $-3\cdot 15\cdot 14\cdot (-4) - 3\cdot 15\cdot 14\cdot (-4)$. No se ve a simple vista pero si separamos en términos podemos notar una particularidad y es que ambos términos son iguales, si lo pensamos a ese termino como $a$, entonces esa operación es equivalente a:
    \begin{gather*}
        -3\cdot 15\cdot 14\cdot (-4) - 3\cdot 15\cdot 14\cdot (-4)=a+a\\
        =2a=2 \cdot (-3)\cdot 15\cdot 14\cdot (-4)
    \end{gather*}
    De allí, trabajando con el teorema fundamental de la aritmetica, se puede reescribir como:
    \begin{gather*}
        2 \cdot (-3)\cdot 15\cdot 14\cdot (-4)= 2  \cdot (-1) \cdot 3 \cdot 3 \cdot 5\cdot 2 \cdot 7 \cdot (-1) \cdot 2^2\\
        =2^4 \cdot (-1)^2 \cdot 3^2 \cdot 5 \cdot 7
    \end{gather*}
    De esta manera podemos trabajar sin necesidad de dividir para saber si es o no múltiplo de los números pedidos pues trabajaríamos directamente con la definición de múltiplos, en este caso como 5 se encuentra entre los factores que forman al número entonces \textbf{es múltiplo de 5} y si descomponemos al $-35$ de igual manera $-35=(-1)\cdot 5 \cdot 7$ esos factores se encuentran en la composición del número que buscamos ver si es múltiplo y por lo tanto resulta también múltiplo de 35.   

    \item[b)] \textit{Encuentren un número entero $k$ de modo que el número $5\cdot(k+3)\cdot 5\cdot(k+3)$ sea múltiplo de $7$. ¿Pueden encontrar otros valores de $k$? ¿Puede haber un valor negativo de $k$?}
    
    \textbf{Resolución:}
    Se busca que $5\cdot(k+3)\cdot 5\cdot(k+3) = 7m$ con $m \in \mathbb{Z}$. Como $\operatorname{mcd}(5, 7) = 1$  debe ser $(k+3)$ múltiplo de $7$, es decir:
    \[
    k + 3 = 7n \iff k = 7n - 3, \quad n \in \mathbb{Z}
    \]
    Para $n=1$, se obtiene $k=4$ ($5^2\cdot(4+3)^2 = 5^2\cdot 7^2$).
    Existen infinitos valores enteros de $k$.
    Sí puede haber valores negativos: por ejemplo para $n=0 \implies k=-3$ ($5^2(-3+3)^2=0=7^2\cdot 0$).
\end{itemize}

\subsubsection{Análisis desde Ponte (2005)}
La primera parte oscila entre un \textbf{ejercicio} y una \textbf{actividad de exploración}, ya que la consigna ``sin hacer cuentas'' demanda analizar la estructura multiplicativa y la factorización en primos más que la ejecución algorítmica que además cuenta con el sentido detrás de esa oración que haya dado la o él docente durante la cursada pues puede haberse trabajado de muchas formas el significado tras hacer cuentas. La segunda parte presenta preguntas guiadas sobre la infinitud y negatividad de las soluciones, situándose como un \textbf{ejercicio con apertura exploratoria} pues, por como debatimos en clases, abre un paraguas amplio sobre que se trabajó anteriormente.

\subsubsection{Análisis desde Skovsmose (2000)}
Corresponde al ámbito de la \textbf{Matemática Pura}. La tarea se ubica en una transición entre el \textbf{Ambiente 1} (Paradigma del Ejercicio) y el \textbf{Ambiente 2} (Escenario de Investigación), al habilitar la búsqueda de patrones y la justificación teórica.

\subsubsection{Diseño Curricular de la Provincia de Córdoba}
\begin{itemize}
    \item \textbf{Curso:} (1. $^{\circ}$ y) 2.$^{\circ}$ Año (Ciclo Básico).
    \item \textbf{Aprendizaje y contenido:} (Uso de las relaciones de múltiplos y divisores de números naturales en situaciones extra e intramatemáticas.)/ Análisis y escritura algebraica de las propiedades de la suma y la multiplicación con números enteros. / Uso de la multiplicación y la división entera con números enteros.
    \item \textbf{Meta:} Operar con números enteros y racionales en situaciones extra e intramatemáticas, y expresar algebraicamente propiedades de las operaciones y relaciones numéricas con números naturales
    \item \textbf{Indicadores de logro:} Comunica mediante expresiones algebraicas las propiedades de la suma y la multiplicación con números enteros. / Anticipa productos y cocientes de números enteros a partir de la comparación e interpretación de los números involucrados.
\end{itemize}
Tomando ese indicador de logro y los temas que lo anteceden como el trabajo con numeros naturales en primer año para anticipar resultados y usar maneras equivalentes de representar números, cómo en este ejercicio se hace el uso de números negativos, está bien ubicarnos en 2do año donde comienzan a trabajar con los enteros, además me sitúo en estos dos contenidos por el hecho de que el primer inciso busca la anticipación de los múltiplos pero el segundo ya busca una expresión algebráica que represente las soluciones posibles
\subsection{Tarea 2}

\subsubsection{Enunciado y Desarrollo Matemático}
La tarea presenta una noticia periodística sobre el promedio de gol por partido en los mundiales de fútbol (1930 a 2010).
\begin{itemize}
    \item[a)] \textit{¿Por qué en Italia 1990 y Sudáfrica 2010 informan el promedio con más de dos cifras decimales?}
    
    \textbf{Resolución:} Porque en Italia 90 el promedio fue $\frac{115}{52} \approx 2{,}2115$ y en Sudáfrica 2010 era $\frac{133}{60} \approx 2{,}2167$. Con dos cifras decimales ambos promedios redondean a $2{,}21$. Se requieren más de dos decimales para poder evidenciar la diferencia de resultados entre ambas divisiones. Aunque no queda del todo claro el uso de puntualmente 4 decimales pues se evidencia diferencia a los 3 y deja eso en manos del escritor del artículo que a una interpretación matemática
    
    \item[b)] \textit{¿Cuáles fueron los tres promedios más altos? ¿Qué observaciones pueden hacer sobre el promedio de goles durante los mundiales?}
    
    \textbf{Resolución:} Los tres más altos fueron: 1) Suiza 1954 ($5{,}38$), 2) Francia 1938 ($4{,}67$) e 3) Italia 1934 ($4{,}12$). Se observa una tendencia histórica decreciente en el promedio de gol a lo largo de las décadas (evolución táctica hacia sistemas más defensivos).
    
    \item[c)] \textit{¿A qué puede deberse que no se jugaron todas las veces la misma cantidad de partidos?}
    
    \textbf{Resolución:} A la modificación del formato del torneo y al incremento paulatino en la cantidad de países participantes clasificados (de 13/16 equipos a 24 y luego a 32 selecciones).
    
    \item[d)] \textit{En Francia 1938 se convirtieron 84 goles en 18 partidos; en Inglaterra 1966 se hicieron 89 goles. ¿Por qué el promedio bajó de 4,67 a 2,78?}
    
    \textbf{Resolución:} Porque el promedio es el cociente $\frac{\text{goles}}{\text{partidos}}$. En 1966 la cantidad de partidos aumentó significativamente (de 18 a 32), mientras que la cantidad de goles apenas aumentó en 5.
\end{itemize}

\subsubsection{Análisis desde Ponte (2005)}
Se clasifica como una \textbf{actividad de exploración} y análisis crítico sobre un texto periodístico. Aunque las preguntas son guiadas, requieren interpretar la media aritmética en contexto e integrar saberes extra-matemáticos y numéricos, lo que sí deja muchos saberes presentados como matemáticos cuando son realmente dados a la interpretación absoluta del lector sobre las intenciones del autor del texto.

\subsubsection{Análisis desde Skovsmose (2000)}
Pertenece al \textbf{Ambiente 5} (Realidad / Paradigma del Ejercicio) con elementos del \textbf{Ambiente 6} (Escenario de Investigación en la Realidad), ya que toma un artículo periodístico real y estimula la reflexión crítica y la modelización del cociente y promedio, además de la comparación entre valores o números dados.

\subsubsection{Diseño Curricular de la Provincia de Córdoba}
\begin{itemize}
    \item \textbf{Curso:} 1.$^{\circ}$ Año (Ciclo Básico).
    \item \textbf{Aprendizaje y contenido:} Uso de la suma, la resta, la multiplicación y la división con números racionales positivos en situaciones extra e intramatemáticas.
    \item \textbf{Meta:} Operar con números enteros y racionales en situaciones extra e intramatemáticas, y expresar algebraicamente propiedades de las operaciones y relaciones numéricas con números naturales
    \item \textbf{Indicadores de logro:} Lee una situación extramatemática, selecciona la información relevante y busca, en distintas fuentes, la información que falta para resolverla.
\end{itemize}
No se hacen calculos muy extensos y se trabaja con los numeros racionales positivos y la división de numeros enteros.
\subsection{Tarea 3}

\subsubsection{Enunciado y Desarrollo Matemático}
\textit{Sea $ABC$ un triángulo rectángulo isósceles cuyos catetos miden $11\text{ cm}$. Entre todos los rectángulos inscriptos en dicho triángulo, ¿hay alguno cuya área sea la mayor de todas? ¿Y la menor? Realizar un gráfico que represente la situación.}

\textbf{Resolución:}
Coloquemos el triángulo en un sistema cartesiano con el vértice recto en $B(0,0)$, el cateto horizontal sobre el eje $x$ hasta $C(11,0)$ y el cateto vertical sobre el eje $y$ hasta $A(0,11)$.
La hipotenusa $AC$ tiene la ecuación de la recta $y = 11 - x$.

Un rectángulo inscripto con un vértice en el origen $(0,0)$ tiene sus lados sobre los ejes y su vértice opuesto $(x, y)$ sobre la hipotenusa $AC$, por lo que $y = 11 - x$, con $x \in (0, 11)$.
El área del rectángulo en función de la base $x$ es:
\[
A(x) = x \cdot y = x(11 - x) = -x^2 + 11x
\]
\begin{itemize}
    \item \textbf{Área máxima:} La función cuadrática $A(x)$ tiene concavidad hacia abajo. El vértice se alcanza en:
    \[
    x_v = -\frac{b}{2a} = \frac{11}{2} = 5{,}5\text{ cm} \implies y = 11 - 5{,}5 = 5{,}5\text{ cm}
    \]
    El área máxima es $A(5{,}5) = 5{,}5 \cdot 5{,}5 = 30{,}25\text{ cm}^2$ (correspondiente a un cuadrado).
    \item \textbf{Área menor:} En el intervalo abierto $(0, 11)$, no existe un valor mínimo absoluto, pues cuando $x \to 0^+$ o $x \to 11^-$, $A(x) \to 0$. Si se consideran rectángulos degenerados en el intervalo cerrado $[0, 11]$, el área mínima es $0\text{ cm}^2$ para $x=0$ o $x=11$.
\end{itemize}

\subsubsection{Análisis desde Ponte (2005)}
Es un \textbf{problema} de optimización geométrica. Presenta un desafío cognitivo alto que requiere vincular la geometría, el álgebra y el estudio de funciones cuadráticas o análisis de simetrías.

\subsubsection{Análisis desde Skovsmose (2000)}
Pertenece al \textbf{Ambiente 1} (Matemática Pura / Paradigma del Ejercicio) o \textbf{Ambiente 2} (Escenario de Investigación en Matemática Pura si se utiliza software de geometría dinámica como GeoGebra para conjeturar el máximo).

\subsubsection{Diseño Curricular de la Provincia de Córdoba}
\begin{itemize}
    \item \textbf{Curso:} 1.$^{\circ}$ Año (Ciclo Básico).
    \item \textbf{Aprendizaje y contenido:} Uso de fórmulas para el cálculo de áreas de cuadriláteros paralelogramos, triángulos y círculos, y perímetros de círculos, en situaciones extra e intramatemáticas.
    \item \textbf{Meta:} Usar propiedades de polígonos convexos, criterios de congruencia y semejanza de triángulos, en situaciones de construcción de figuras geométricas y cálculo de áreas y volúmenes.
    \item \textbf{Indicadores de logro:} Elabora, explica y compara fórmulas propias y de otras/os que permiten calcular el área de cuadriláteros, paralelogramos y triángulos.
\end{itemize}
    Si bien trabaja profundamente con geometría y las relaciones entre figuras, la mayoría del desarrollo tiene un análisis algebráico y del estudio de las funciones de fondo que amplía la dificultad y diría que no sólo se camufla en la parte geométrica pura sino que modeliza situaciones de variaciones lineales queriendo dar el paso a el análisis de funciones cuadráticas y su representación
\subsection{Tarea 4}

\subsubsection{Enunciado y Desarrollo Matemático}
\textit{Dados dos ángulos $A$ y $B$. Construir, si es posible, un triángulo que tenga un ángulo igual a $A$ y otro ángulo igual a $B$. ¿Es posible construir dos distintos? ¿Por qué? ¿Será cierto que dados dos ángulos, siempre es posible construir un triángulo?}

\textbf{Resolución:}
\begin{enumerate}
    \item \textbf{Condición de posibilidad:} En la geometría euclidiana plana, la suma de los ángulos interiores de un triángulo es $180^\circ$. Por ende, es posible construir un triángulo si y sólo si:
    \[
    \hat{A} + \hat{B} < 180^\circ
    \]
    El tercer ángulo queda unívocamente determinado por $\hat{C} = 180^\circ - (\hat{A} + \hat{B})$. Si $\hat{A} + \hat{B} \ge 180^\circ$, es imposible construirlo.
    \item \textbf{Construcción de triángulos distintos:} Sí, es posible construir infinitos triángulos no congruentes pero semejantes (criterio de semejanza Ángulo-Ángulo, AA). Al fijar únicamente dos ángulos, los triángulos difieren por una homotecia (escalamiento de la longitud de sus lados).
    \item \textbf{Validez general:} No es cierto que siempre sea posible; falla cuando la suma de los dos ángulos dados es mayor o igual a dos rectos ($180^\circ$).
\end{enumerate}

\subsubsection{Análisis desde Ponte (2005)}
Es una \textbf{actividad de exploración} y conjeturación geométrica. Exige a los estudiantes poner a prueba propiedades de existencia y unicidad de figuras geométricas, así como diferenciar congruencia de semejanza.

\subsubsection{Análisis desde Skovsmose (2000)}
Se encuadra en el \textbf{Ambiente 2} (Matemática Pura / Escenario de Investigación), ya que promueve la exploración geométrica constructiva y el planteo de hipótesis mediante regla y compás o software dinámico.

\subsubsection{Diseño Curricular de la Provincia de Córdoba}
\begin{itemize}
    \item \textbf{Curso:} 1.$^{\circ}$ Año (Ciclo Básico).
    \item \textbf{Aprendizaje y contenido:} Formulación de argumentaciones sobre condiciones geométricas de triángulos y cuadriláteros convexos y propiedades de pares de ángulos, en situaciones de construcción y comunicación.
    \item \textbf{Meta:} Usar propiedades de polígonos convexos, criterios de congruencia y semejanza de triángulos, en situaciones de construcción de figuras geométricas y cálculo de áreas y volúmenes.
    \item \textbf{Indicadores de logro:} Fundamenta la cantidad de triángulos diferentes que se pueden construir a partir de la información sobre sus lados y ángulos.
\end{itemize}

Se Trabaja mucho con triángulos y conocimiento básico sobre sus ángulos pero hay una profundización en el trabajo sobre la unicidad que les permite dar los primeros pasos a prueba y error que los indicadores de logro plantean correctamente en esta sección.
\subsection{Tarea 5}

\subsubsection{Enunciado y Desarrollo Matemático}
\textit{Considerar rectángulos de perímetro $P = 24\text{ cm}$.}
\begin{itemize}
    \item[a)] Si $2b + 2h = 24 \implies b + h = 12 \implies h = 12 - b$. Para valores de base $b \in \{2, 4, 6, 8, 10\}$, las alturas son $h \in \{10, 8, 6, 4, 2\}$.
    \item[b)] En términos geométricos reales (lados positivos), $0 < b < 12$. El límite inferior es $0$ y el superior es $12$ (no alcanzables si el rectángulo no es degenerado).
    \item[c)] Los puntos $(b, h)$ están alineados sobre la recta $h = -b + 12$.
    \item[d)] Función altura: $h(b) = 12 - b$.
    \item[e)] El área $A = b \cdot h = b(12 - b)$ no es constante; varía parabólicamente.
    \item[f-h)] Función cuadrática de área: $A(b) = -b^2 + 12b$.
    \item[g)] Vértice en $b = 6\text{ cm}$, con $h = 6\text{ cm}$ (cuadrado), alcanzando un área máxima de $A_{\max} = 36\text{ cm}^2$.
    \item[i)] Dominios contextualizados: $\operatorname{Dom}(h) = (0, 12)$ y $\operatorname{Dom}(A) = (0, 12)$.
\end{itemize}

\subsubsection{Análisis desde Ponte (2005)}
Es una secuencia articulada estructurada como una \textbf{actividad de exploración guiada} que transita desde un nivel de \textbf{ejercicio} (tabulación inicial) hacia un \textbf{problema} de modelización funcional y optimización.

\subsubsection{Análisis desde Skovsmose (2000)}
Ubicada en el \textbf{Ambiente 1} (Matemática Pura / Paradigma del Ejercicio) con fuerte andamiaje hacia el \textbf{Ambiente 2}, al estudiar el pasaje entre registros de representación semiótica (tabla, gráfico cartesiano, fórmula algebraica).

\subsubsection{Diseño Curricular de la Provincia de Córdoba}
\begin{itemize}
    \item \textbf{Curso:} 4.$^{\circ}$ Año (Ciclo Orientado).
    \item \textbf{Aprendizaje y contenido:} Análisis de funciones lineales f(x) = kx + b, presentadas en textos, tablas, gráficos y fórmulas, en situaciones extra e intramatemáticas.
    \item \textbf{Meta:} Modelizar situaciones vinculadas a la función lineal, cuadrática, racional, exponencial, logarítmica y trigonométricas.
    \item \textbf{Indicadores de logro:} Todos los presentes en esta meta y año
\end{itemize}

Se busca que en una situación de geometría se pueda manejar y determinar el uso de información y el uso de funciones y su profundo análisis sobre el comportamiento y dominio.
\subsection{Tarea 6}

\subsubsection{Enunciado y Desarrollo Matemático}
\textit{Llenado de recipientes y gráfica de la altura $h$ en función del volumen $V$.}
\begin{itemize}
    \item[a)] \textbf{Recipiente cilíndrico:} La tabla muestra una tasa de variación constante: $\frac{h}{V} = \frac{1}{10} = 0{,}1\text{ cm/cm}^3$. La gráfica $h(V) = 0{,}1 V$ es una recta que pasa por el origen (función lineal de proporcionalidad directa).
    \item[b)] \textbf{Otros recipientes:}
    \begin{itemize}
        \item \textit{Botella con cuello angosto:} Sección ancha abajo (pendiente baja constante) y sección estrecha arriba (pendiente alta constante). Gráfica con cambio brusco de pendiente positiva.
        \item \textit{Embudo (cono invertido sobre cilindro pequeño):} Base cilíndrica (recta empinada) y luego sección transversal cónica que se ensancha con la altura, por lo que la tasa de llenado disminuye (función cóncava hacia abajo).
        \item \textit{Recipiente esférico con cuello:} Sección se ensancha hasta el ecuador (pendiente decreciente) y luego se angosta (pendiente creciente), formando un punto de inflexión; finaliza con un tramo lineal empinado para el cuello cilíndrico.
        \item \textit{Cono truncado:} Sección disminuye uniformemente con la altura, por lo que la altura sube cada vez más rápido por cada unidad de volumen (función convexa, cóncava hacia arriba).
    \end{itemize}
\end{itemize}

\subsubsection{Análisis desde Ponte (2005)}
Constituye una \textbf{actividad de exploración} y análisis cualitativo de gráficos funcionales. Promueve el razonamiento sobre el comportamiento del volumen sin depender exclusivamente de fórmulas algebraicas explícitas.

\subsubsection{Análisis desde Skovsmose (2000)}
Se clasifica en el \textbf{Ambiente 3} (Semirrealidad / Paradigma del Ejercicio) o \textbf{Ambiente 4} (Semirrealidad / Escenario de Investigación), ya que utiliza una situación física idealizada para construir nociones de cálculo y covariación.

\subsubsection{Diseño Curricular de la Provincia de Córdoba}
\begin{itemize}
    \item \textbf{Curso:} 2.$^{\circ}$ Año (Ciclo Básico).
    \item \textbf{Aprendizaje y contenido:} Uso de fórmulas para el cálculo de volúmenes de prismas, pirámides y cuerpos redondos, en situaciones extra e intramatemáticas.
    \item \textbf{Meta:} Usar propiedades de polígonos convexos, criterios de congruencia y semejanza de triángulos, en situaciones de construcción de figuras geométricas y cálculo de áreas y volúmenes.
    \item \textbf{Indicadores de logro:} Selecciona las fórmulas de acuerdo con la información disponible en la situación leída, para el cálculo del volumen de cuerpos compuestos por prismas, pirámides y/o cuerpos redondos, y justifica la elección realizada.
\end{itemize}
Si bien no hacen el uso directo de las fórmulas si hacen un análisis de su conocimiento en el volumen de distintos cuerpos y su relación con el análisis de las funciones que representan y su comportamiento en el cambio de forma del cuerpo. Lo que me conflictua a elegir si están situados en esta parte del curriculum o es posterior al estudio del comportamiento de las funciones.
\subsection{Tarea 8}

\subsubsection{Enunciado y Desarrollo Matemático}
\textit{¿Cuál es el perímetro de un cuadrado cuya área es de $23\text{ cm}^2$?}

\textbf{Resolución:}
El área de un cuadrado de lado $l$ es $A = l^2$.
Dado que $A = 23\text{ cm}^2$:
\[
l = \sqrt{23}\text{ cm}
\]
El perímetro $P$ es cuatro veces la longitud del lado:
\[
P = 4 \cdot l = 4\sqrt{23}\text{ cm} \approx 4 \cdot 4{,}7958 = 19{,}183\text{ cm}
\]

\subsubsection{Análisis desde Ponte (2005)}
Es un típico \textbf{ejercicio} cerrado y de desafío cognitivo reducido. Consiste en la aplicación directa y sucesiva de dos fórmulas elementales de geometría plana y radicación.

\subsubsection{Análisis desde Skovsmose (2000)}
Pertenece claramente al \textbf{Ambiente 1} (Matemática Pura / Paradigma del Ejercicio), caracterizado por enunciados estándar de libro de texto con una única respuesta exacta y técnica operatoria prefijada.

\subsubsection{Diseño Curricular de la Provincia de Córdoba}
\begin{itemize}
    \item \textbf{Curso:} 1.$^{\circ}$ Año (Ciclo Básico).
    \item \textbf{Aprendizaje y contenido:} Uso de fórmulas para el cálculo de áreas de cuadriláteros paralelogramos, triángulos y círculos, y perímetros de círculos, en situaciones extra e intramatemáticas.
    \item \textbf{Meta:} Usar propiedades de polígonos convexos, criterios de congruencia y semejanza de triángulos, en situaciones de construcción de figuras geométricas y cálculo de áreas y volúmenes.
    \item \textbf{Indicadores de logro:} Selecciona las fórmulas de acuerdo a la información disponible en la situación leída, para el cálculo del área de figuras compuestas por paralelogramos, triángulos y círculos, y justifica su elección.
\end{itemize}
Es una actividad bastante situada que responde al indicador de logro del contenido elegido.
\subsection{Tarea 9}

\subsubsection{Enunciado y Desarrollo Matemático}
La tarea trabaja sobre el Cuadro H17 del Censo 2010 del INDEC acerca de hogares con población originaria por provincia.
\begin{itemize}
    \item[a)] \textbf{Variables analizadas:}
    \begin{itemize}
        \item \textit{Provincia / Jurisdicción:} Variable cualitativa nominal.
        \item \textit{Total de hogares:} Variable cuantitativa discreta.
        \item \textit{Hogares con personas indígenas u originarias:} Variable cuantitativa discreta.
        \item \textit{Proporción o porcentaje relativo:} Se calcula como $\frac{\text{Hogares Indígenas}}{\text{Total de Hogares}} \cdot 100$.
    \end{itemize}
    \item[b)] \textbf{Observaciones relevantes:} La provincia de Buenos Aires concentra la mayor cantidad absoluta de hogares indígenas ($121.385$), pero en términos relativos, provincias del NOA y Patagonia (como Jujuy con $\approx 11{,}1\%$, Chubut con $\approx 11{,}2\%$, Río Negro con $\approx 9{,}3\%$) presentan porcentajes sensiblemente superiores al promedio nacional ($\approx 3{,}03\%$).
    \item[c)] \textbf{Visualización y Estadística:} Gráficos de barras comparativos (totales vs.\ porcentajes) y cálculo de medidas de posición (media, mediana, cuartiles, valores extremos).
\end{itemize}

\subsubsection{Análisis desde Ponte (2005)}
Es una \textbf{actividad de investigación} escolar en el ámbito de la estadística cívica y social. La consigna es abierta, demanda recolección bibliográfica, formulación de hipótesis, procesamiento de datos y redacción de conclusiones socio-demográficas fundamentadas.

\subsubsection{Análisis desde Skovsmose (2000)}
Pertenece al \textbf{Ambiente 6} (Realidad / Escenario de Investigación). Los estudiantes abordan datos censales verídicos, confrontando la dimensión sociopolítica con las herramientas estadísticas.

\subsubsection{Diseño Curricular de la Provincia de Córdoba}
\begin{itemize}
    \item \textbf{Curso:} 3.$^{\circ}$ Año (Ciclo Básico).
    \item \textbf{Aprendizaje y contenido:} Análisis conjunto de datos recolectados de dos variables estadísticas cualitativas y cuantitativas discretas.
    \item \textbf{Meta:} Cuantificar la incertidumbre recurriendo a las nociones de probabilidad clásica y comunicar información estadística.
    \item \textbf{Indicadores de logro:} Todos los del contenido.
\end{itemize}
Trabajan directamente con el objetivo de ejrcitar y reafianzar los indicadores de logro del contenido como "Justifica la elección del gráfico estadístico más apropiado para presentar la información conjunta de dos variables cualitativas y/o cuantitativas discretas" y su organización dentro de esos aspectos.
\end{document}
